
\documentclass{article}
\usepackage{colortbl}
\usepackage{makecell}
\usepackage{multirow}
\usepackage{supertabular}

\begin{document}

\newcounter{utterance}

\centering \large Interaction Transcript for game `taboo', experiment `medium\_en', episode 15 with reflexion{-}llama{-}3.1{-}8b{-}instruct.
\vspace{24pt}

{ \footnotesize  \setcounter{utterance}{1}
\setlength{\tabcolsep}{0pt}
\begin{supertabular}{c@{$\;$}|p{.15\linewidth}@{}p{.15\linewidth}p{.15\linewidth}p{.15\linewidth}p{.15\linewidth}p{.15\linewidth}}
    \# & $\;$A & \multicolumn{4}{c}{Game Master} & $\;\:$B\\
    \hline

    \theutterance \stepcounter{utterance}  
    & & \multicolumn{4}{p{0.6\linewidth}}{
        \cellcolor[rgb]{0.9,0.9,0.9}{
            \makecell[{{p{\linewidth}}}]{
                \texttt{\tiny{[P1$\langle$GM]}}
                \texttt{You have attempted this same task before and did not succeed. These are the plans you wrote after those attempts.} \\
\texttt{Trial \#1: New plan: I failed to provide a description that would lead the other player to guess the target word within three attempts. In my last clue, I gave away too much information by saying "You sit on it," which directly relates to the target word "toilet." Next time, I will provide a more general description that still hints at the target word without giving away its identity. For example, I could say "It's a fixture in a bathroom where you can spend a few minutes." This description still relates to the target word but doesn't give away its name.} \\
\\ 
\texttt{Trial \#2: New plan: In my last attempt, I failed to provide a description that would lead the other player to guess the target word within three attempts. I gave away too much information by saying "It's often cleaned daily," which is a common trait of the target word "toilet" but also applies to other related words like "restroom." Next time, I will provide a more specific description that still hints at the target word without giving away its identity. For example, I could say "It's a fixture in a bathroom where you can spend a few minutes, and it's often associated with a specific sound." This description still relates to the target word but doesn't give away its name.} \\
\\ 
\\ 
\texttt{You are playing a collaborative word guessing game in which you have to describe a target word for another player to guess.} \\
\\ 
\texttt{Rules:} \\
\texttt{(a) You have to reply in the form: CLUE: <some text>. Guesses from the other player will start with GUESS.} \\
\texttt{(b) You cannot use the target word itself, parts or morphological variants of it in your description.} \\
\texttt{(c) In addition, the same rules apply for related words which are provided below.} \\
\\ 
\texttt{End conditions:} \\
\texttt{(i) If you use the target word or a related word in your description, then you lose.} \\
\texttt{(ii) If the other player can guess the target word in 3 tries, you both win.} \\
\\ 
\texttt{Let us start.} \\
\\ 
\texttt{This is the target word that you need to describe and that the other player needs to guess:} \\
\\ 
\texttt{toilet} \\
\\ 
\texttt{Related words are:} \\
\\ 
\texttt{{-} flush} \\
\texttt{{-} restroom} \\
\texttt{{-} urinate} \\
\\ 
\texttt{Important: You are under time pressure, give short descriptions that are to the point!} \\
            }
        }
    }
    & \\ \\

    \theutterance \stepcounter{utterance}  
    & \multicolumn{4}{p{0.6\linewidth}}{
        \cellcolor[rgb]{0.9,0.9,0.9}{
            \makecell[{{p{\linewidth}}}]{
                \texttt{\tiny{[P1$\rangle$GM]}}
                \texttt{CLUE: It's a fixture in a private space where people often spend a few minutes.} \\
            }
        }
    }
    & & \\ \\

    \theutterance \stepcounter{utterance}  
    & & & \multicolumn{2}{p{0.3\linewidth}}{
        \cellcolor[rgb]{0.9,0.9,0.9}{
            \makecell[{{p{\linewidth}}}]{
                \texttt{\tiny{[GM$|$GM]}}
                \texttt{continue} \\
            }
        }
    }
    & & \\ \\

    \theutterance \stepcounter{utterance}  
    & & & \multicolumn{2}{p{0.3\linewidth}}{
        \cellcolor[rgb]{0.9,0.9,0.9}{
            \makecell[{{p{\linewidth}}}]{
                \texttt{\tiny{[GM$|$GM]}}
                \texttt{its a fixture in a private space where people often spend a few minutes} \\
            }
        }
    }
    & & \\ \\

    \theutterance \stepcounter{utterance}  
    & & \multicolumn{4}{p{0.6\linewidth}}{
        \cellcolor[rgb]{0.9,0.9,0.9}{
            \makecell[{{p{\linewidth}}}]{
                \texttt{\tiny{[GM$\rangle$P2]}}
                \texttt{You have attempted this same task before and did not succeed. These are the plans you wrote after those attempts.} \\
\texttt{Trial \#1: New plan: } \\
\\ 
\texttt{In my last attempt, I was able to guess the first two clues correctly, but failed to guess the target word in the third trial. Upon reflection, I realize that I was too focused on guessing a specific word based on the given clues, rather than trying to understand the relationship between the clues and the target word. } \\
\\ 
\texttt{In the first trial, the clue was "A room in a house where you go to do a private activity." I correctly guessed the word "BATHROOM". } \\
\\ 
\texttt{In the second trial, the clue was "It's often found inside a bathroom." I correctly guessed the word "SHOWER". } \\
\\ 
\texttt{In the third trial, the clue was "You sit on it." I incorrectly guessed the word "TOILET". } \\
\\ 
\texttt{Upon reflection, I realize that I should have paid closer attention to the relationship between the clues. The first clue mentioned a room, the second clue mentioned something found inside a bathroom, and the third clue mentioned something you sit on. } \\
\\ 
\texttt{A more effective strategy would be to try to understand the relationship between the clues and use that information to make a more informed guess. For example, I could have guessed a word that is often found inside a bathroom and is something you sit on, such as a "STOOL". } \\
\\ 
\texttt{Next time, I will pay closer attention to the relationship between the clues and use that information to make a more informed guess.} \\
\\ 
\texttt{Trial \#2: New plan: } \\
\\ 
\texttt{In my last attempt, I was able to guess the first two clues correctly, but failed to guess the target word in the third trial. Upon reflection, I realize that I was too focused on guessing a specific word based on the given clues, rather than trying to understand the relationship between the clues and the target word. } \\
\\ 
\texttt{In the first trial, the clue was "A place where you can take care of personal hygiene." I correctly guessed the word "BATHROOM". } \\
\\ 
\texttt{In the second trial, the clue was "A common activity done there takes only a few seconds." I correctly guessed the word "SINK". } \\
\\ 
\texttt{In the third trial, the clue was "It's often cleaned daily." I incorrectly guessed the word "SINK" again. } \\
\\ 
\texttt{Upon reflection, I realize that I should have paid closer attention to the relationship between the clues. The first clue mentioned a place for personal hygiene, the second clue mentioned a common activity that takes only a few seconds, and the third clue mentioned something that is often cleaned daily. } \\
\\ 
\texttt{A more effective strategy would be to try to understand the relationship between the clues and use that information to make a more informed guess. For example, I could have guessed a word that is often found in a bathroom and is cleaned daily, such as a "MIRROR". } \\
\\ 
\texttt{Next time, I will pay closer attention to the relationship between the clues and use that information to make a more informed guess. I will also try to think of words that fit the description of all the clues, rather than just one or two.} \\
\\ 
\\ 
\texttt{You are playing a collaborative word guessing game in which you have to guess a target word that another player describes to you.} \\
\\ 
\texttt{You can make one guess at each trial. You win when you guess the target word. You lose when you cannot guess it in 3 tries.} \\
\\ 
\texttt{After each trial you will get a new hint from the other player which starts with CLUE.} \\
\\ 
\texttt{Make your guesses by just saying the word using the following form: GUESS: <a word>} \\
\\ 
\texttt{Let us start.} \\
\\ 
\\ 
\texttt{CLUE: It's a fixture in a private space where people often spend a few minutes.} \\
            }
        }
    }
    & \\ \\

    \theutterance \stepcounter{utterance}  
    & & & \multicolumn{4}{p{0.6\linewidth}}{
        \cellcolor[rgb]{0.9,0.9,0.9}{
            \makecell[{{p{\linewidth}}}]{
                \texttt{\tiny{[GM$\langle$P2]}}
                \texttt{GUESS: CHAIR} \\
            }
        }
    }
     \\ \\

    \theutterance \stepcounter{utterance}  
    & & & \multicolumn{2}{p{0.3\linewidth}}{
        \cellcolor[rgb]{0.9,0.9,0.9}{
            \makecell[{{p{\linewidth}}}]{
                \texttt{\tiny{[GM$|$GM]}}
                \texttt{continue} \\
            }
        }
    }
    & & \\ \\

    \theutterance \stepcounter{utterance}  
    & & & \multicolumn{2}{p{0.3\linewidth}}{
        \cellcolor[rgb]{0.9,0.9,0.9}{
            \makecell[{{p{\linewidth}}}]{
                \texttt{\tiny{[GM$|$GM]}}
                \texttt{chair} \\
            }
        }
    }
    & & \\ \\

    \theutterance \stepcounter{utterance}  
    & & \multicolumn{4}{p{0.6\linewidth}}{
        \cellcolor[rgb]{0.9,0.9,0.9}{
            \makecell[{{p{\linewidth}}}]{
                \texttt{\tiny{[P1$\langle$GM]}}
                \texttt{GUESS: CHAIR} \\
            }
        }
    }
    & \\ \\

    \theutterance \stepcounter{utterance}  
    & \multicolumn{4}{p{0.6\linewidth}}{
        \cellcolor[rgb]{0.9,0.9,0.9}{
            \makecell[{{p{\linewidth}}}]{
                \texttt{\tiny{[P1$\rangle$GM]}}
                \texttt{CLUE: It's typically found in a room with a sink and a mirror.} \\
            }
        }
    }
    & & \\ \\

    \theutterance \stepcounter{utterance}  
    & & & \multicolumn{2}{p{0.3\linewidth}}{
        \cellcolor[rgb]{0.9,0.9,0.9}{
            \makecell[{{p{\linewidth}}}]{
                \texttt{\tiny{[GM$|$GM]}}
                \texttt{continue} \\
            }
        }
    }
    & & \\ \\

    \theutterance \stepcounter{utterance}  
    & & & \multicolumn{2}{p{0.3\linewidth}}{
        \cellcolor[rgb]{0.9,0.9,0.9}{
            \makecell[{{p{\linewidth}}}]{
                \texttt{\tiny{[GM$|$GM]}}
                \texttt{its typically found in a room with a sink and a mirror} \\
            }
        }
    }
    & & \\ \\

    \theutterance \stepcounter{utterance}  
    & & \multicolumn{4}{p{0.6\linewidth}}{
        \cellcolor[rgb]{0.9,0.9,0.9}{
            \makecell[{{p{\linewidth}}}]{
                \texttt{\tiny{[GM$\rangle$P2]}}
                \texttt{CLUE: It's typically found in a room with a sink and a mirror.} \\
            }
        }
    }
    & \\ \\

    \theutterance \stepcounter{utterance}  
    & & & \multicolumn{4}{p{0.6\linewidth}}{
        \cellcolor[rgb]{0.9,0.9,0.9}{
            \makecell[{{p{\linewidth}}}]{
                \texttt{\tiny{[GM$\langle$P2]}}
                \texttt{GUESS: VANITY} \\
            }
        }
    }
     \\ \\

    \theutterance \stepcounter{utterance}  
    & & & \multicolumn{2}{p{0.3\linewidth}}{
        \cellcolor[rgb]{0.9,0.9,0.9}{
            \makecell[{{p{\linewidth}}}]{
                \texttt{\tiny{[GM$|$GM]}}
                \texttt{continue} \\
            }
        }
    }
    & & \\ \\

    \theutterance \stepcounter{utterance}  
    & & & \multicolumn{2}{p{0.3\linewidth}}{
        \cellcolor[rgb]{0.9,0.9,0.9}{
            \makecell[{{p{\linewidth}}}]{
                \texttt{\tiny{[GM$|$GM]}}
                \texttt{vanity} \\
            }
        }
    }
    & & \\ \\

    \theutterance \stepcounter{utterance}  
    & & \multicolumn{4}{p{0.6\linewidth}}{
        \cellcolor[rgb]{0.9,0.9,0.9}{
            \makecell[{{p{\linewidth}}}]{
                \texttt{\tiny{[P1$\langle$GM]}}
                \texttt{GUESS: VANITY} \\
            }
        }
    }
    & \\ \\

    \theutterance \stepcounter{utterance}  
    & \multicolumn{4}{p{0.6\linewidth}}{
        \cellcolor[rgb]{0.9,0.9,0.9}{
            \makecell[{{p{\linewidth}}}]{
                \texttt{\tiny{[P1$\rangle$GM]}}
                \texttt{CLUE: It's often associated with a specific sound when it's used.} \\
            }
        }
    }
    & & \\ \\

    \theutterance \stepcounter{utterance}  
    & & & \multicolumn{2}{p{0.3\linewidth}}{
        \cellcolor[rgb]{0.9,0.9,0.9}{
            \makecell[{{p{\linewidth}}}]{
                \texttt{\tiny{[GM$|$GM]}}
                \texttt{continue} \\
            }
        }
    }
    & & \\ \\

    \theutterance \stepcounter{utterance}  
    & & & \multicolumn{2}{p{0.3\linewidth}}{
        \cellcolor[rgb]{0.9,0.9,0.9}{
            \makecell[{{p{\linewidth}}}]{
                \texttt{\tiny{[GM$|$GM]}}
                \texttt{its often associated with a specific sound when its used} \\
            }
        }
    }
    & & \\ \\

    \theutterance \stepcounter{utterance}  
    & & \multicolumn{4}{p{0.6\linewidth}}{
        \cellcolor[rgb]{0.9,0.9,0.9}{
            \makecell[{{p{\linewidth}}}]{
                \texttt{\tiny{[GM$\rangle$P2]}}
                \texttt{CLUE: It's often associated with a specific sound when it's used.} \\
            }
        }
    }
    & \\ \\

    \theutterance \stepcounter{utterance}  
    & & & \multicolumn{4}{p{0.6\linewidth}}{
        \cellcolor[rgb]{0.9,0.9,0.9}{
            \makecell[{{p{\linewidth}}}]{
                \texttt{\tiny{[GM$\langle$P2]}}
                \texttt{GUESS: SHOWER} \\
            }
        }
    }
     \\ \\

    \theutterance \stepcounter{utterance}  
    & & & \multicolumn{2}{p{0.3\linewidth}}{
        \cellcolor[rgb]{0.9,0.9,0.9}{
            \makecell[{{p{\linewidth}}}]{
                \texttt{\tiny{[GM$|$GM]}}
                \texttt{continue} \\
            }
        }
    }
    & & \\ \\

    \theutterance \stepcounter{utterance}  
    & & & \multicolumn{2}{p{0.3\linewidth}}{
        \cellcolor[rgb]{0.9,0.9,0.9}{
            \makecell[{{p{\linewidth}}}]{
                \texttt{\tiny{[GM$|$GM]}}
                \texttt{shower} \\
            }
        }
    }
    & & \\ \\

    \theutterance \stepcounter{utterance}  
    & & & \multicolumn{2}{p{0.3\linewidth}}{
        \cellcolor[rgb]{0.9,0.9,0.9}{
            \makecell[{{p{\linewidth}}}]{
                \texttt{\tiny{[GM$|$GM]}}
                \texttt{3} \\
            }
        }
    }
    & & \\ \\

    \theutterance \stepcounter{utterance}  
    & \multicolumn{4}{p{0.6\linewidth}}{
        \cellcolor[rgb]{0.95,0.95,0.95}{
            \makecell[{{p{\linewidth}}}]{
                \texttt{\tiny{[P2$\rightarrow$P1]}}
                \texttt{Post{-}episode Reflexion (not sent to GM):} \\
\texttt{New plan: In my last attempts, I failed to provide a description that would lead the other player to guess the target word within three attempts. I gave away too much information by providing descriptions that were too general or too specific. Next time, I will provide a description that is more balanced, giving the player a general idea of the target word's context and function, without revealing its identity. I will also try to introduce new information in each clue that is not too obvious or too vague, to keep the player guessing. For example, I could say "It's a fixture in a bathroom where you can spend a few minutes, and it's often associated with a specific sound, but it's not something you can sit on." This description still relates to the target word but doesn't give away its name, and introduces a new piece of information that the player can use to make their next guess.} \\
            }
        }
    }
    & & \\ \\

    \theutterance \stepcounter{utterance}  
    & & & \multicolumn{4}{p{0.6\linewidth}}{
        \cellcolor[rgb]{0.95,0.95,0.95}{
            \makecell[{{p{\linewidth}}}]{
                \texttt{\tiny{[P1$\rightarrow$P2]}}
                \texttt{Post{-}episode Reflexion (not sent to GM):} \\
\texttt{New plan: } \\
\\ 
\texttt{In my last attempt, I was able to guess the first two clues correctly, but failed to guess the target word in the third trial. Upon reflection, I realize that I was too focused on guessing a specific word based on the given clues, rather than trying to understand the relationship between the clues and the target word. } \\
\\ 
\texttt{In the first trial, the clue was "A fixture in a private space where people often spend a few minutes." I correctly guessed the word "CHAIR". } \\
\\ 
\texttt{In the second trial, the clue was "It's typically found in a room with a sink and a mirror." I correctly guessed the word "VANITY". } \\
\\ 
\texttt{In the third trial, the clue was "It's often associated with a specific sound when it's used." I incorrectly guessed the word "SHOWER". } \\
\\ 
\texttt{Upon reflection, I realize that I should have paid closer attention to the relationship between the clues. The first clue mentioned a fixture in a private space, the second clue mentioned a room with a sink and a mirror, and the third clue mentioned a specific sound associated with its use. } \\
\\ 
\texttt{A more effective strategy would be to try to understand the relationship between the clues and use that information to make a more informed guess. For example, I could have guessed a word that is often found in a room with a sink and a mirror, and is associated with a specific sound, such as a "WASHBASIN". } \\
\\ 
\texttt{Next time, I will pay closer attention to the relationship between the clues and use that information to make a more informed guess. I will also try to think of words that fit the description of all the clues, rather than just one or two.} \\
            }
        }
    }
     \\ \\

\end{supertabular}
}

\end{document}
